% !PRINTED 624
\weberpage{624}
The group $\Psi$ is isomorphic with a group, contained in $\mathfrak{N}$ (\S.\,169), of integral
rational numbers taken with respect to the modulus $m$, which we denote by $\mathfrak{A}$, and
which consists of all numbers $a$ satisfying a
congruence
\weq{3}{a \equiv p^h \pmod{m}}
and which we can therefore represent symbolically by
\weq{4}{\mathfrak{A} \equiv p^h \pmod{m}}
if $h$ denotes an arbitrary integral exponent (cf. \S.\,18).

If, then, we select the numbers $\xi_1, \xi_2, \ldots, \xi_e$ from $\mathfrak{N}$
suitably, we can set:
\weq{5}{\mathfrak{N} = \mathfrak{A}\xi_1 + \mathfrak{A}\xi_2 + \mathfrak{A}\xi_3 + \cdots + \mathfrak{A}\xi}
% [sic] source prints the final \xi without the expected subscript e; (6)--(7) supply the context
and to each of these subsidiary groups there corresponds one of the conjugate functionals $\mathfrak{p}_{\xi_1}, \mathfrak{p}_{\xi_2}, \ldots, \mathfrak{p}_{\xi_e}$, which are all different from one another and whose product is associated with $p$.
% [sic] "Functionale" as printed; context elsewhere on the page uses "Primfactoren"/"Primideale" for this notion

We therefore set
\weq{6}{p = \mathfrak{p}_{\xi_1} \mathfrak{p}_{\xi_2} \ldots \mathfrak{p}_{\xi_e}.}
But in order to characterise more exactly the system of numbers
\weq{7}{\xi_1, \xi_2, \ldots, \xi_e,}
on which the matter chiefly depends,
we remark that for this purpose we can take any system of $e$ numbers of the
group $\mathfrak{N}$, provided only that no two among them,
$\xi, \xi'$, satisfy a congruence
\[
\xi'\xi^{-1} \equiv p^h \pmod{m}.
\]

In order to find such a system of numbers, it is well to separate the case
of an odd $m$ from the other, in which $m$ is a
power of 2.

1. If first $m = q^{\varkappa}$ is odd, we take a
primitive root $g$ of $m$ (\S.\,14) and set
\weq{8}{p \equiv g^{\gamma} \pmod{m}.}
For $f$ is [by (1), (2)] the smallest positive number that satisfies the
condition
\[
\gamma f \equiv 0 \pmod{\mu}
\]
and consequently $e$ is the greatest common divisor
of $\mu$ and $\gamma$, and any number $n \equiv g^{\lambda}$ is only then
%CONT>>
